%% USE one of these:
%% * the first when using pdflatex, which directly typesets your document in the
%%   chosen pdf/a format and you want to publish your thesis online,

%% * the second when you want to print your thesis to bind it, or
%% * the third when producing a ps file and a pdf/a from it.
%%
\documentclass[english, 12pt, a4paper, sci, utf8, a-2b, online]{aaltothesis}
%\documentclass[english, 12pt, a4paper, elec, utf8, a-1b]{aaltothesis}
%\documentclass[english, 12pt, a4paper, elec, dvips, online]{aaltothesis}

%% Use the following options in the \documentclass macro above:
%% your school: arts, biz, chem, elec, eng, sci
%% the character encoding scheme used by your editor: utf8, latin1
%% thesis language: english, finnish, swedish
%% make an archiveable PDF/A-1b or PDF/A-2b compliant file: a-1b, a-2b
%%                    (with pdflatex, a normal pdf containing metadata is
%%                     produced without the a-*b option)
%% typeset in symmetric layout and blue hypertext for online publication: online
%%            (no option is the default, resulting in a wide margin on the
%%             binding side of the page and black hypertext)
%% two-sided printing: twoside (default is one-sided printing)
%%

%% Use one of these if you write in Finnish (see the Finnish template
%% opinnaytepohja.tex)
%\documentclass[finnish, 12pt, a4paper, elec, utf8, a-1b, online]{aaltothesis}
%\documentclass[finnish, 12pt, a4paper, elec, utf8, a-1b]{aaltothesis}
%\documentclass[finnish, 12pt, a4paper, elec, dvips, online]{aaltothesis}

\usepackage{graphicx}

%% Math fonts, symbols, and formatting; these are usually needed
\usepackage{amsfonts,amssymb,amsbsy,amsmath}

%% Change the school field to specify your school if the automatically set name
%% is wrong
% \university{aalto-yliopisto}
% \school{Sähkötekniikan korkeakoulu}

%% Edit to conform to your degree programme
%%
\degreeprogram{Computer, Communication and Information Sciences}
%%

%% Your major
%%
\major{Machine Learning, Data Science and Artificial Intelligence}
%%

%% Major subject code
%%
\code{SCI3044}
%%
 
%% Choose one of the three below
%%
% \univdegree{BSc}
\univdegree{MSc}
%\univdegree{Lic}
%%

%% Your name (self explanatory...)
%%
\thesisauthor{Joel Himanen}
%%

%% Your thesis title comes here and possibly again together with the Finnish or
%% Swedish abstract. Do not hyphenate the title, and avoid writing too long a
%% title. Should LaTeX typeset a long title unsatisfactorily, you mght have to
%% force a linebreak using the \\ control characters.
%% In this case...
%% Remember, the title should not be hyphenated!
%% A possible "and" in the title should not be the last word in the line, it
%% begins the next line.
%% Specify the title again without the linebreak characters in the optional
%% argument in box brackets. This is done because the title is part of the 
%% metadata in the pdf/a file, and the metadata cannot contain linebreaks.
%%
\thesistitle{Title TBD}
%\thesistitle[Title of the thesis]{Title of\\ the thesis}
%%

%%
\place{Espoo}
%%

%% The date for the bachelor's thesis is the day it is presented
%%
\date{31.12.2022}
%%

%% Thesis supervisor
%% Note the "\" character in the title after the period and before the space
%% and the following character string.
%% This is because the period is not the end of a sentence after which a
%% slightly longer space follows, but what is desired is a regular interword
%% space.
%%
\supervisor{Prof.\ Alexander Jung}
%%

%% Advisor(s)---two at the most---of the thesis. Check with your supervisor how
%% many official advisors you can have.
%%
\advisor{D.Sc. (Tech.) Sami Jokela}
\advisor{Prof. Marko Turpeinen}
%%

%% Aaltologo: syntax:
%% \uselogo{aaltoRed|aaltoBlue|aaltoYellow|aaltoGray|aaltoGrayScale}{?|!|''}
%% The logo language is set to be the same as the thesis language.
%%
\uselogo{aaltoBlue}{''}
%%

%% The English abstract:
%% All the details (name, title, etc.) on the abstract page appear as specified
%% above.
%% Thesis keywords:
%% Note! The keywords are separated using the \spc macro
%%
\keywords{Insert\spc Keywords\spc Here}
%%

%% The abstract text. This text is included in the metadata of the pdf file as well
%% as the abstract page.
%%
\thesisabstract{
Your abstract in English. Keep the abstract short. The abstract explains your 
research topic, the methods you have used, and the results you obtained. In the 
PDF/A format of this thesis, in addition to the abstract page, the abstract text is 
written into the pdf file's metadata. Write here the text that goes into the 
metadata. The metadata cannot contain special characters, linebreak or paragraph 
break characters, so these must not be used here. If your abstract does not contain 
special characters and it does not require paragraphs, you may take advantage of 
the abstracttext macro (see the comment below). Otherwise, the metadata abstract 
text must be identical to the text on the abstract page.
}

%% Copyright text. Copyright of a work is with the creator/author of the work
%% regardless of whether the copyright mark is explicitly in the work or not.
%% You may, if you wish, publish your work under a Creative Commons license (see
%% creaticecommons.org), in which case the license text must be visible in the
%% work. Write here the copyright text you want. It is written into the metadata
%% of the pdf file as well.
%% Syntax:
%% \copyrigthtext{metadata text}{text visible on the page}
%% 
%% In the macro below, the text written in the metadata must have a \noexpand
%% macro before the \copyright special character, and macros (\copyright and
%% \year here) must be separated by the \ character (space chacter) from the
%% text that follows. The macros in the argument of the \copyrighttext macro
%% automatically insert the year and the author's name. (Note! \ThesisAuthor is
%% an internal macro of the aaltothesis.cls class file).
%% Of course, the same text could have simply been written as
%% \copyrighttext{Copyright \noexpand\copyright\ 2018 Eddie Engineer}
%% {Copyright \copyright{} 2018 Eddie Engineer}
%%
\copyrighttext{Copyright \noexpand\copyright\ \number\year\ \ThesisAuthor}
{Copyright \copyright{} \number\year{} \ThesisAuthor}

%% You can prevent LaTeX from writing into the xmpdata file (it contains all the 
%% metadata to be written into the pdf file) by setting the writexmpdata switch
%% to 'false'. This allows you to write the metadata in the correct format
%% directly into the file thesistemplate.xmpdata.
%\setboolean{writexmpdatafile}{false}

%% All that is printed on paper starts here
%%
\begin{document}

%% Create the coverpage
%%
\makecoverpage

%% Typeset the copyright text.
%% If you wish, you may leave out the copyright text from the human-readable
%% page of the pdf file. This may seem like a attractive idea for the printed
%% document especially if "Copyright (c) yyyy Eddie Engineer" is the only text
%% on the page. However, the recommendation is to print this copyright text.
%%
\makecopyrightpage

%% Note that when writting your thesis in English, place the English abstract
%% first followed by the possible Finnish or Swedish abstract.

%% Abstract text
%% All the details (name, title, etc.) on the abstract page appear as specified
%% above.
%%
\begin{abstractpage}[english]
  Your abstract in English. Keep the abstract short. The abstract explains your
  research topic, the methods you have used, and the results you obtained.

  The abstract text of this thesis is written on the readable abstract page as
  well as into the pdf file's metadata via the $\backslash$thesisabstract macro
  (see above). Write here the text that goes onto the readable abstract page.
  You can have special characters, linebreaks, and paragraphs here. Otherwise,
  this abstract text must be identical to the metadata abstract text.

  If your abstract does not contain special characters and it does not require
  paragraphs, you may take advantage of the abstracttext macro (see the comment
  below).
\end{abstractpage}

%% The text in the \thesisabstract macro is stored in the macro \abstractext, so
%% you can use the text metadata abstract directly as follows:
%%
%\begin{abstractpage}[english]
%	\abstracttext{}
%\end{abstractpage}

%% Force a new page so that the possible Finnish or Swedish abstract does not
%% begin on the same page
%%
\newpage
%%
%% Abstract in Finnish.  Delete if you don't need it. 
%%
\thesistitle{Otsikko suomeksi}
\supervisor{Prof.\ Alexander Jung}
\advisor{TkT Sami Jokela}
\advisor{Prof. Marko Turpeinen}
\degreeprogram{Selvitä}
%\department{Elektroniikan ja nanotekniikan laitos}
\major{Selvitä}
%% The keywords need not be separated by \spc now.
\keywords{avainsana1, avainasana2}
%% Abstract text
\begin{abstractpage}[finnish]
  Tiivistelmässä on lyhyt selvitys
  kirjoituksen tärkeimmästä sisällöstä: mitä ja miten on tutkittu,
  sekä mitä tuloksia on saatu.
\end{abstractpage}


%% Preface
%%
%% This section is optional. Remove it if you do not want a preface.
\mysection{Preface}
%\mysection{Esipuhe}
Preface.

\vspace{5cm}
Otaniemi, 31.12.2022

\vspace{5mm}
{\hfill Joel Himanen \hspace{1cm}}

%% Force a new page after the preface
%%
\newpage


%% Table of contents. 
%%
\thesistableofcontents


%% Symbols and abbreviations
\mysection{Abbreviations}

\begin{tabular}{ll}
  \textbf{API} & application programming interface \\
  \textbf{CE}  & circular economy                  \\
  \textbf{GBC} & Green Building Council            \\
  \textbf{GHG} & greenhouse gas                    \\
\end{tabular}

%% \clearpage is similar to \newpage, but it also flushes the floats (figures
%% and tables).
%%
\cleardoublepage

%% Text body begins. Note that since the text body is mostly in Finnish the
%% majority of comments are also in Finnish after this point. There is no point
%% in explaining Finnish-language specific thesis conventions in English.
%% This text will be translated to English soon.
%%

\mysection{Acknowledgements}

%% Leave page number of the first page empty
\thispagestyle{empty}

\clearpage

\section{Introduction} \label{intro}

% Contribution: 5\%
% Order of work: 6/6

%% In a thesis, every section starts a new page, hence \clearpage
\clearpage

\section{Background} \label{background}

% Contribution: 25\%
% Order of work: 1/6

This chapter establishes the reasoning behind the urgent need for a global transition towards a circular economy (CE) and describes the role digitalization plays in the process.
It takes the Finnish construction industry as an example of this transition, and specifies the key challenges and opportunities of driving the change with data and digital solutions.

\subsection{Circular economy and digitalization} \label{ce}

It is no news that the ongoing environmental crises, such as global warming and loss of biodiversity, pose a large part of the most
significant threats of our lifetime. The challenge of steering humankind towards a sustainable way of life is monumental, and requires
extraordinay feats, as well as cross-disciplinary coordination and co-operation from all entities ranging from individual people to
businesses, industries, countries and international alliances.

In 2020, the Finnish Innovation Fund Sitra conducted a study on the megatrends for the beginning decade~\cite{dufva}, identifying
five different, but intertwined global megatrends, three of which are the following:
\begin{itemize}
  \item \textit{The urgent need for ecological reconstruction}
  \item \textit{The redefinition of the economy}
  \item \textit{Technology is embedded in everything}
\end{itemize}
Getting the interactions between technology and economy to work in favour of fighting environmental crises is an incredibly complex, yet critical
challenge. In other words, it is safe to state that the two latter trends play a key role in the way we drive ecological reconstruction in the years to come.

The culprits behind our environmental crises are manifold, but at the end of the day, much of the problem boils
down to how we use the natural resources our planet provides us with. Already in 2019,
resource extraction and processing made up approximately half of global greenhouse gas (GHG) emissions, in addition to more than 90\% of
biodiversity loss and water stress.~\cite{oberle} It has become evident, that the sun has set on the old linear economy model. We
can not afford to extract, use and dispose of resources with the pace we have reached due to assumptions of limitless economic growth
with little to no regard to external costs. The concept of a circular economy presents an alternative model to this, where product
lifetime and reusability is maximized, production waste is minimized, and business models based on sharing and reusing products are implemented.
The potential in these strategies is enormous: according to a Material Economics study from 2018~\cite{enkvist}, the transition to a more
circular economy could reduce GHG emissions from heavy industry in the EU by 56\% by 2050. Contrary to a common misconception, this shift does not
need to happen at the cost of economic growth as a whole. Aligning technologies and business models to tackle wasteful value creation chains
and shift towards a more circular model would actually create new markets, and boost employment, lifting EU's GDP by 7\% and
yielding annual benefits of up to €1,8 trillion by 2030.~\cite{macarthur}

As Sitra states, technology is embedded in everything, and the transition to CE is, of course, no exception. The breathtaking pace digitalization
has developed with during the last decades has presented us with countless possibilities and opportunities, but also with its own set of challenges.
For example, with increasing amounts of digital solutions and services increases also the demand for energy. Hence, it would be naïve to assume
digitalization to be a "silver bullet" that directly solves the problems related to implementing CE models, as stated in a 2019 study by the European
Policy Centre (EPC).~\cite{hedberg} That being said, the study focuses on establishing the role of digitalization as a key enabler of CE -
as long as it is properly managed. They divide approaches to use digitalization to accelerate the transition to CE into three categories:
\begin{enumerate}
  \item Improving knowledge, connections and information sharing
  \item Making business models, products and processes more circular
  \item Strengthening the roles of citizens and consumers
\end{enumerate}
Within the scope of this thesis, we will mainly focus on category 1, exploring the ways proper data aggregation, utilization and sharing can support businesses in adapting more circular
models. Specifically, we will look at the construction industry, as it is responsible for almost half of global raw material extraction and a quarter of
global GHG emissions.~\cite{amro}

\subsection{Circularity in the Finnish construction industry} \label{finnish_construction}

When reviewing the environmental impact of the construction industry, it is essential to divide carbon emissions into three categories according to the phases
in the lifecycle of a building:
\begin{enumerate}
  \item \textbf{Upfront carbon} (materials production and construction phase)
  \item \textbf{Operational carbon} (use phase, energy consumption)
  \item \textbf{End of life carbon} (demolition phase, demolition waste handling)
\end{enumerate}
These are defined by the World Green Building Council (WorldGBC) in their 2019 report "Bringing embodied carbon upfront"~\cite{adams}. The same study finds that with the
recent advances in the fied of renewable energy, operational carbon emissions are reduced significantly. Consequently, upfront and end of life carbon will
grow quickly in proportion. In particular, the WorldGBC predicts that upfront carbon will be responsible for up to half of the whole carbon footprint of new construction between the
the present and 2050. A similar trend can be seen in Finland, as well. Up until 2020, approximately 80\% of the Finnish built environment's total lifecycle emissions came from operational
carbon.~\cite{laine} With low-carbon energy gaining foothold in the Finnish energy market, focus is shifting increasingly towards reducing upfront and end of life carbon by transitioning
to a more circular operating model. The Finnish government has recently taken the lead in driving this transition across industries in Finland. In the spring of 2021,
they established a strategic programme \cite{gov} to pave the way towards a new Finnish economy based on CE by 2035. The programme dedicates a significant amount of focus on the
construction industry, and specifies several action points to promote CE in its context, including financial incentives, digitalization and data management.

There are, however, plenty of challenges to overcome. Especially crucial is the disconnection between supply and demand in the Finnish recycled construction product market. On a demolition site,
the prime contractor owns all of the materials (although the responsibility of demolition waste handling lies with the demolition contractor, if the demolition is outsourced by the prime contractor),
and therefore is also responsible for processing all of the material streams released during demolition accordingly, including all related costs and taxes. On the other hand, there are plenty of costs
embedded in the process of preparing the materials from extraction to reuse, and said costs often outweigh the costs of disposing of these materials, which makes it less appealing for contractors
to salvage them for reuse (some components can be extracted pre-demolition and used again as such) and recycling (products that are not feasible for reuse are processed into new products).
It can thus be argued, that there exists at least a potential financial incentive for the contractors to sell as much of the released materials as possible,
to be used in new construction. The challenge is to minimize the costs of processing used construction products and demolition waste for reuse and recycling. The problem described above also
has the effect of making virgin materials often more attractive price-wise, which hinders the demand for used and recycled construction products, and makes adapting
a more circular model difficult from the point of view of procurement as well.

This problematic is clearly a key obstacle, as it effects negatively both the supply and demand sides of the non-virgin construction product market.
This has also been recognized by the Finnish government, as it has recently set in motion investigation projects to e.g. identify
ambiguities in construction product qualification procedures.~\cite{zhu} As an example result, the Finnish Ministry of the Environment released a policy brief~\cite{tahtinen} on verifying the
feasability of reused materials in June 2022. It states that it is sufficient to check for feasabilty with construction site -specific verification, i.e. a CE-marking is not
required for construction products that are reused as such.

However, even if there was a strong financial incentive to favor non-virgin materials, a significant challenge remains aligning demand
with supply on a temporal level. Currently, contractors are able to release information of available materials only as soon as they become available, as the demolition project proceeds.
Because storage possibilities for these materials are few and often costly, the time window to find a buyer is far too narrow. To counter this challenge, Finnish officials are working
on a series of databases and registers that would together provide a centralized data storage system for information on demolition projects and plans, construction product
charecteristics and material streams, to name a few. These systems are planned to come with diverse APIs to enable private entities to develop their own services and platforms to utilize all of
the data stored in these public systems. The following section will introduce one such project, in which the author is also involved through his employment.

\subsection{Motiva's data platform and ecosystem for material circulations} \label{motiva}

In order to make the raw data from the existing and upcoming registers for construction materials and demolition waste described in the previous section available and usable to companies and
consumers, it is vital that different services are developed to refine and process the data. Even though work on the state-wide registers is still in progress, the Finnish state-owned
company Motiva Ltd. is already working on a platform for construction material circulation. Motiva is a company of experts specialized in providing organizations and consumers
with information, solutions and services on how to use energy and materials more efficiently and sustainably. With this platform, it is Motiva's aim to build a hub that collects
information from various sources under a single platform, and acts as a meeting place for supply and demand of non-virgin construction materials.

At the moment, this platform is still in the planning phase. Motiva has developed the concept for it, and it is currently being validated and further developed in two different demolition projects.
The first one takes place in Vattuniemi, Helsinki, where 16 buildings are set to be demolished during 2022-23. The other demolition site is located in the Tampere area, where especially the reuse
and circulation of ashes as a construction product is being explored. These projects are vital in gathering field experience on the challenges and requirements of material circulation.
Specifically, Motiva seeks to pilot planned features of their platform, and to gain knowledge on how exactly predictive, data-driven demolition planning supports a circular economy
in the construction industry.

\subsection{The potential of pre-demolition audits} \label{pre_dem}

A critical part of shifting towards predictive demolition planning are the procedures that guide the mapping of the demolition site. These include mapping of harmful, but also reusable
and recyclable materials that the building to be demolished consists of. Historically, only the mapping of asbestos and other harmful materials has been required by law, but nowadays it is increasingly
common to do a broader material mapping. The most extensive mapping process is the pre-demolition audit (in Finnish "purkukartoitus"), which also includes, for example, reusable movables and
recommendations on how to to reuse, recycle or dispose of the material streams released during demolition.~\cite{lehtonen} These audits are at the moment voluntary, but officials are constantly planning ways to
encourage real estate owners and contractors to use them as a tool to increase material circulation. This is slightly challenging, since a more extensive audit is also more expensive to the client.
It is, however, crucial from the point of view of CE, that pre-demolition audits are carried out as carefully as possible, because gaining knowledge of future material streams is necessary in
order to align demand with supply in the non-virgin material market.

Carrying out pre-demolition audits, and especially estimating the amounts of different materials released, is hard to do accurately. All experts interviewed for this thesis agreed on this,
and stated several possible reasons for the significant errors that occur even in state-of-the-art audits. Some of the material amounts are especially hard to predict, because they might
often be hidden from plain sight in the foundations of the building in question. An example of this kind of material is mineral wool, which is often used as an insulator. Especially in the case of old
buildings, where the quality of material documentation might be suboptimal, these kind of "hidden" materials can be nigh unto impossible to predict. It might also happen, that some renovations or
other changes during the building's lifetime have not been properly documented, which causes some material streams to come as a surprise during demolition. Finally, it is also natural that clients
want to minimize the costs of a demolition project, and the pre-demolition auditors may have to compromise extensiveness and accuracy of the audit in order to stay within the budget.

Given the potential of pre-demolition audits from a CE point of view, and the nature of the challenges described above, it is evident that making the audit process more data-driven holds significant
possibilities to develop construction material circulation. The obvious bottleneck at the moment, however, is the lack of data. There are no standardized, mandatory procedures to store pre-demolition
audit, material mapping, or realized material stream data. This leads to the relatively small amount of feasible data to be scattered all over different stakeholders, and to have little to
no quality guarantees. This challenge is at least partly being addressed by the projects described in sections \ref{finnish_construction} and \ref{motiva}, but this thesis will contribute
to the ways of data aggregation on a more detailed level, and explore the possibilities of utilizing this data in the future.

\clearpage

\section{Model and methods}

% Contribution: 20\%
% Order of work: 2/6

\subsection{Data}
\subsection{Features}
\subsection{Prediction targets}
\subsection{Modeling approach and algorithms}

\clearpage

\section{Results}

% Contribution: 10\%
% Order of work: 3/6

\subsection{Methods and results of model testing}


\clearpage

\section{Discussion}

% Contribution: 35\%
% Order of work: 4/6

\subsection{Meaning of modeling results}
\subsection{Roadmap for data aggregation and preservation}
\subsection{Future modeling possibilities}

\clearpage

\section{Summary}

% Contribution: 5\%
% Order of work: 5/6


%% Handle bibliography with bib(la)tex, commented out
\clearpage

\bibliographystyle{plain}
\bibliography{refs}

%\thesisbibliography

%\begin{thebibliography}{99}

%\end{thebibliography}

%% Appendices
%% If you don't have appendices, remove \clearpage and \thesisappendix below.
% \clearpage

% \thesisappendix

% \section{Example appendix}



\end{document}
